% !TEX root = ../main.tex
% !TEX program = lualatex
\documentclass[../main.tex]{subfiles}
\IfSubfilesClassLoaded{\externaldocument{\subfix{../build/main}}}{}
% ====================
% File: 17B_Risk_Management_and_Problem_Solving.tex
% ====================
\begin{document}
\IfSubfilesClassLoaded{\chapter{Risk Management \& Problem Solving}}{}
% === Risk Management & Problem Solving (EDO 3.3.6) ===
\section{Risk Management \& Problem Solving}

\subsection{Summary}
\begin{description}
  \item[\textbf{Definition.}] Acquisition risk is the inability to achieve cost, schedule, and performance objectives; risk is defined by probability, consequence, and root cause~\autocite{edo-3-3-6-risk-management-problem-solving-2025,DoD-RIO-Guide-2017}.
  \item[\textbf{Sources.}] Risks arise from external drivers (threat, funding, contractor, politics) and internal drivers (technology, design, manufacturing, support, management)~\autocite{edo-3-3-6-risk-management-problem-solving-2025}.
  \item[\textbf{APB linkage.}] The \ac{apb} sets the baseline used to plan, control, and manage risk across cost, schedule, and performance goals~\autocite{edo-3-3-6-risk-management-problem-solving-2025}.
  \item[\textbf{Process.}] Risk management is a continuous cycle: plan, identify, analyze, handle, and monitor; opportunity and issue management use related but distinct handling choices~\autocite{edo-3-3-6-risk-management-problem-solving-2025,DoD-RIO-Guide-2017}.
  \item[\textbf{Problem solving.}] Trade studies, \ac{maut}, fishbone diagrams, the 5 Whys, and multi-voting support objective decisions and root-cause discovery~\autocite{edo-3-3-6-risk-management-problem-solving-2025}.
\end{description}

\subsection{What Is Acquisition Risk?}
\begin{description}
  \item[\textbf{Risk management.}] Identifies, evaluates, and prioritizes risks, then applies resources to minimize, monitor, and control probability or impact~\autocite{edo-3-3-6-risk-management-problem-solving-2025}.
  \item[\textbf{Program management.}] Risk management is integral to program management and systems engineering~\autocite{edo-3-3-6-risk-management-problem-solving-2025}.
  \item[\textbf{Risk components.}] Probability of occurrence, consequence (impact), and root cause define each risk~\autocite{edo-3-3-6-risk-management-problem-solving-2025,DoD-RIO-Guide-2017}.
\end{description}

\subsection{General Sources of Acquisition Risk}
\begin{description}
  \item[\textbf{External.}] Threat and requirements, funding, contractor capability, and politics~\autocite{edo-3-3-6-risk-management-problem-solving-2025}.
  \item[\textbf{Internal.}] Technology, design/engineering, manufacturing, support/logistics, cost/schedule, and management~\autocite{edo-3-3-6-risk-management-problem-solving-2025}.
\end{description}

\subsection{APB and Risk Management}
\begin{description}
  \item[\textbf{Baseline.}] The \ac{apb} documents cost, schedule, and performance thresholds/objectives agreed by the \ac{pm} and \ac{mdauthority}~\autocite{edo-3-3-6-risk-management-problem-solving-2025}.
  \item[\textbf{Planning anchor.}] The \ac{apb} provides the baseline used to plan, control, and manage risk to attain program goals~\autocite{edo-3-3-6-risk-management-problem-solving-2025}.
  \item[\textbf{Ground rules.}] Analyze risks with realistic timeframes, consider event timing, and identify risks at the lowest feasible \ac{wbs} level~\autocite{edo-3-3-6-risk-management-problem-solving-2025}.
\end{description}

\subsection{Risk Management Model}
\begin{description}
  \item[\textbf{Planning.}] Define the program approach, roles, resources, and battle rhythm; document in \ac{sep} and Acquisition Strategy; often combined in a \ac{rio} plan or \ac{romb}~\autocite{edo-3-3-6-risk-management-problem-solving-2025}.
  \item[\textbf{Identification.}] Determine what can go wrong and document root causes using program data, test results, and trends~\autocite{edo-3-3-6-risk-management-problem-solving-2025}.
  \item[\textbf{Analysis.}] Size risk by probability and consequence, use risk matrices, and isolate root causes (the 5 Whys)~\autocite{edo-3-3-6-risk-management-problem-solving-2025}.
  \item[\textbf{Handling.}] Define responses, resources, ownership, and reporting requirements~\autocite{edo-3-3-6-risk-management-problem-solving-2025}.
  \item[\textbf{Monitoring.}] Track burn-down plans, update risk status, and adjust actions based on performance metrics~\autocite{edo-3-3-6-risk-management-problem-solving-2025}.
\end{description}

\subsection{Risk Identification Inputs}
\begin{itemize}
  \item Current/proposed staffing, resources and dependencies, test failures, negative data trends, and performance shortfalls~\autocite{edo-3-3-6-risk-management-problem-solving-2025}.
  \item \ac{tpm}, life-cycle cost information, \ac{wbs}, \ac{ims}, \ac{evm}, and technical baseline maturity~\autocite{edo-3-3-6-risk-management-problem-solving-2025}.
\end{itemize}

\subsection{Risk Mitigation Techniques}
\begin{itemize}
  \item Prototyping, modeling and simulation, and technology demonstrations with decision points~\autocite{edo-3-3-6-risk-management-problem-solving-2025}.
  \item Intelligence analyses, data dependencies, and threat projections~\autocite{edo-3-3-6-risk-management-problem-solving-2025}.
  \item Multiple design approaches, alternative designs, and phased activities to address high-risk areas early~\autocite{edo-3-3-6-risk-management-problem-solving-2025}.
  \item Manufacturability, industrial base availability, supply-chain analysis, and independent risk assessments~\autocite{edo-3-3-6-risk-management-problem-solving-2025}.
  \item Schedule and funding margins sized to identified risks~\autocite{edo-3-3-6-risk-management-problem-solving-2025}.
\end{itemize}

\subsection{Risk Handling Options}
\begin{description}
  \item[\textbf{Control.}] Reduce probability/impact (\ac{p3i}, reuse software, parallel design, additional testing)~\autocite{edo-3-3-6-risk-management-problem-solving-2025}.
  \item[\textbf{Avoid.}] Change course to eliminate the risk (redesign, reduce scope, extend schedule, use \ac{cots})~\autocite{edo-3-3-6-risk-management-problem-solving-2025}.
  \item[\textbf{Accept.}] Continue the plan while identifying resources or schedule to respond to consequences~\autocite{edo-3-3-6-risk-management-problem-solving-2025}.
  \item[\textbf{Transfer.}] Reassign or delegate mitigation tasks, and sometimes financial exposure, to another entity; transfer does not eliminate Government responsibility, and schedule/performance risk cannot be fully transferred because the Government still needs the capability~\autocite{edo-3-3-6-risk-management-problem-solving-2025,DoD-RIO-Guide-2017}.
\end{description}

\subsection{Risk Monitoring}
\begin{description}
  \item[\textbf{Status and communication.}] Share risk status with stakeholders and review burn-down plans through periodic \ac{ipt} meetings and \acp{ibr}~\autocite{edo-3-3-6-risk-management-problem-solving-2025}.
  \item[\textbf{Tracking.}] Update the risk matrix, alert leadership when plans need adjustment, and evaluate impacts across milestones~\autocite{edo-3-3-6-risk-management-problem-solving-2025}.
\end{description}

\subsection{Risk Management Terminology}
\begin{description}
  \item[\textbf{Concern.}] A newly identified item before \ac{rmwg} review~\autocite{edo-3-3-6-risk-management-problem-solving-2025}.
  \item[\textbf{Candidate.}] A concern accepted for analysis or a segment risk elevated to program level and presented to the \ac{ermb}~\autocite{edo-3-3-6-risk-management-problem-solving-2025}.
  \item[\textbf{Active.}] Risk managed under a formal handling plan in planning or tracking sub-phases~\autocite{edo-3-3-6-risk-management-problem-solving-2025}.
  \item[\textbf{Watch.}] Low-level risk monitored without a handling plan; elevated if risk increases~\autocite{edo-3-3-6-risk-management-problem-solving-2025}.
  \item[\textbf{Retired/Rejected.}] Retired risks are mitigated or no longer relevant; rejected risks are removed due to low concern~\autocite{edo-3-3-6-risk-management-problem-solving-2025}.
  \item[\textbf{Issue.}] A realized risk or problem whose consequence has occurred or is certain~\autocite{edo-3-3-6-risk-management-problem-solving-2025}.
\end{description}

\subsection{Risk Management Software Tools}
\begin{description}
  \item[\textbf{Standards.}] Tools often align to Project Management Institute, National Institute of Standards and Technology, or International Standards Organization guidance~\autocite{edo-3-3-6-risk-management-problem-solving-2025}.
  \item[\textbf{Capabilities.}] Risk registers, modeling/forecasting, Monte Carlo analysis, Office compatibility, and web access~\autocite{edo-3-3-6-risk-management-problem-solving-2025}.
\end{description}

\subsection{Opportunity Management}
\begin{description}
  \item[\textbf{Definition.}] Opportunities are future uncertainties that can exceed performance, accelerate schedules, or reduce cost~\autocite{edo-3-3-6-risk-management-problem-solving-2025}.
  \item[\textbf{Parallel process.}] Uses the same organization and process as risk management; current \ac{dod} opportunity options are pursue now, defer, reevaluate, or reject~\autocite{edo-3-3-6-risk-management-problem-solving-2025,DoD-RIO-Guide-2017}.
\end{description}

\subsection{Issue Management}
\begin{description}
  \item[\textbf{Definition.}] Issues are realized risks or problems that already occurred or are certain to occur~\autocite{edo-3-3-6-risk-management-problem-solving-2025}.
  \item[\textbf{Handling.}] Similar to risk management but without likelihood assessment because probability is 1. Primary issue options are ignore, meaning accept consequences based on business-case analysis, or control, meaning implement corrective action; avoid and transfer are less common issue options~\autocite{edo-3-3-6-risk-management-problem-solving-2025,DoD-RIO-Guide-2017}.
\end{description}

\subsection{Trade Study Fundamentals}
\begin{description}
  \item[\textbf{Purpose.}] Evaluate viable alternatives and justify the preferred choice for requirements, design, or processes~\autocite{edo-3-3-6-risk-management-problem-solving-2025}.
  \item[\textbf{Core trade.}] Always includes analysis of cost versus performance~\autocite{edo-3-3-6-risk-management-problem-solving-2025}.
\end{description}

\subsection{Generic Problem Solving Model}
\begin{enumerate}
  \item Define the problem using facts, flowcharts, and cause-effect diagrams~\autocite{edo-3-3-6-risk-management-problem-solving-2025}.
  \item Generate alternatives through \ac{ipt} brainstorming~\autocite{edo-3-3-6-risk-management-problem-solving-2025}.
  \item Evaluate and select an alternative~\autocite{edo-3-3-6-risk-management-problem-solving-2025}.
  \item Implement the solution and follow up~\autocite{edo-3-3-6-risk-management-problem-solving-2025}.
\end{enumerate}

\subsection{Cause and Effect (Fishbone) Diagram}
\begin{description}
  \item[\textbf{Purpose.}] Identifies root causes using an Ishikawa diagram~\autocite{edo-3-3-6-risk-management-problem-solving-2025}.
  \item[\textbf{Six cause categories.}] Materials, machines, methods, measurement, manpower, and milieu (environment)~\autocite{edo-3-3-6-risk-management-problem-solving-2025}.
\end{description}

\subsection{The 5 Whys}
\begin{description}
  \item[\textbf{Method.}] Iterative questioning to reach the underlying root cause; the number five is a guideline~\autocite{edo-3-3-6-risk-management-problem-solving-2025}.
  \item[\textbf{Best practices.}] Use an accurate problem statement, unbiased answers, and focus on root causes not symptoms~\autocite{edo-3-3-6-risk-management-problem-solving-2025}.
\end{description}

\subsection{Multi-Voting}
\begin{description}
  \item[\textbf{Purpose.}] Narrows long lists of options using group judgment and weighted voting~\autocite{edo-3-3-6-risk-management-problem-solving-2025}.
  \item[\textbf{How.}] Select top items individually, tally votes (weighted if needed), and repeat until priorities are clear~\autocite{edo-3-3-6-risk-management-problem-solving-2025}.
\end{description}

\subsection{Multi-Attribute Utility Theory (MAUT)}
\begin{enumerate}
  \item Identify alternatives~\autocite{edo-3-3-6-risk-management-problem-solving-2025}.
  \item Choose attributes and define a numerical grading scale~\autocite{edo-3-3-6-risk-management-problem-solving-2025}.
  \item Apply weights to attributes~\autocite{edo-3-3-6-risk-management-problem-solving-2025}.
  \item Score each alternative on each attribute~\autocite{edo-3-3-6-risk-management-problem-solving-2025}.
  \item Multiply scores by weights and sum~\autocite{edo-3-3-6-risk-management-problem-solving-2025}.
  \item Compare results~\autocite{edo-3-3-6-risk-management-problem-solving-2025}.
\end{enumerate}

\subsection{Quick Review}
\begin{itemize}
  \item Risk is comprised of which three variables~\autocite{edo-3-3-6-risk-management-problem-solving-2025}? \textbf{Answer:} Root cause, probability of occurrence, and consequence.
  \item Risk mitigation option described as lowering event probability through multiple tests is which~\autocite{edo-3-3-6-risk-management-problem-solving-2025}? \textbf{Answer:} Control.
  \item What is the issue-management trap~\autocite{DoD-RIO-Guide-2017}? \textbf{Answer:} Do not score likelihood; an issue has occurred or is certain, so focus on consequence and corrective action.
  \item ``Use of proven technologies in commercial and non-developmental items eliminates technical risk''~\autocite{edo-3-3-6-risk-management-problem-solving-2025}? \textbf{Answer:} False.
\end{itemize}

%====================
% End of file
%====================
\IfSubfilesClassLoaded{
  \RenewDocumentCommand{\entryname}{}{\textbf{\color{Modern} Acronym}}
  \RenewDocumentCommand{\descriptionname}{}{\textbf{\color{Modern} Definition}}
  \printnoidxglossary[
    type=\acronymtype,
    title=Acronyms,
    style=long-booktabs
  ]}{}
\end{document}
